% Part: propositional-logic
% Chapter: syntax-and-semantics
% Section: formulas

\documentclass[../../../include/open-logic-section]{subfiles}

\begin{document}

\olfileid[pt]{pl}{syn}{fml}

\olsection{\usetoken{P}{formula} proposicionais}

As !!^{formula}s da lógica proposicional são construídas a partir de
\emph{!!{propositional
    variable}s}\iftag{prvFalse}{\iftag{prvTrue}{,}{ e} da constante
  proposicional~$\lfalse$}{}\iftag{prvTrue}{\iftag{prvFalse}{
    e}{} da constante proposicional~$\ltrue$}{} usando \emph{conectivos
  lógicos}.

\begin{enumerate}
\item !!^a{denumerable} conjunto~$\PVar$ de !!{propositional variable}s $\Obj p_0$,
  $\Obj p_1$, \dots
\tagitem{prvFalse}{A constante proposicional para !!{falsity}~$\lfalse$.}{}
\tagitem{prvTrue}{A constante proposicional para !!{truth}~$\ltrue$.}{}
\item Os conectivos lógicos:
  \startycommalist
  \iftag{prvNot}{\ycomma $\lnot$ (negação)}{}%
  \iftag{prvAnd}{\ycomma $\land$ (conjunção)}{}%
  \iftag{prvOr}{\ycomma $\lor$ (disjunção)}{}%
  \iftag{prvIf}{\ycomma $\lif$ (!!{conditional})}{}%
  \iftag{prvIff}{\ycomma $\liff$ (!!{biconditional})}{}%
\item Sinais de pontuação: (, ) e a vírgula.
\end{enumerate}

Denotamos essa linguagem da lógica proposicional por $\Lang L_0$.

\iftag{defNot,defOr,defAnd,defIf,defIff,defTrue,defFalse,defEx,defAll}{%
Além dos conectivos primitivos introduzidos acima, também usamos os
seguintes símbolos \emph{definidos}:
\startycommalist
  \iftag{defNot}{\ycomma $\lnot$ (negação)}{}%
  \iftag{defAnd}{\ycomma $\land$ (conjunção)}{}%
  \iftag{defOr}{\ycomma $\lor$ (disjunção)}{}%
  \iftag{defIf}{\ycomma $\lif$ (!!{conditional})}{}%
  \iftag{defIff}{\ycomma $\liff$ (!!{biconditional})}{}%
  \iftag{defFalse}{\ycomma $\lfalse$ (!!{falsity})}{}%
  \iftag{defTrue}{\ycomma $\ltrue$ (!!{truth})}}{}.

\begin{tagblock}{defNot,defOr,defAnd,defIf,defIff,defTrue,defFalse,defEx,defAll}
\begin{explain}
Um símbolo definido não faz oficialmente parte da linguagem, mas é
introduzido como abreviação informal: ele nos permite abreviar
fórmulas que ficariam bastante longas se usássemos apenas símbolos
primitivos. Essa é uma vantagem evidente. A vantagem maior, porém, é
que as demonstrações ficam mais curtas. Se um símbolo é primitivo, ele
precisa ser tratado separadamente nas demonstrações. Portanto, quanto
mais símbolos primitivos houver, mais longas serão nossas
demonstrações.
\end{explain}
\end{tagblock}

% Alternate symbols

\begin{intro}
Talvez você conheça terminologia e símbolos diferentes dos que usamos
acima. Textos de lógica (e professores) costumam usar $\sim$, $\neg$
e~!\ para a ``negação'', e $\wedge$, $\cdot$ e $\&$ para a
``conjunção''. Símbolos comumente usados para a ``condicional'' ou a
``implicação'' são $\rightarrow$, $\Rightarrow$ e $\supset$.
\iftag{prvIff,defIff}{Os símbolos para a ``bicondicional'', a
  ``bi-implicação'' ou a ``equivalência (material)'' são $\leftrightarrow$,
  $\Leftrightarrow$ e $\equiv$.}{}
\iftag{prvFalse,defFalse}{O símbolo $\lfalse$ recebe diversos nomes:
  ``falsidade'', ``falso'', ``absurdo'' ou ``base''.}{}
\iftag{prvTrue,defTrue}{O símbolo $\ltrue$ recebe diversos nomes:
  ``verdade'', ``verdadeiro'' ou ``topo''.}{}
\end{intro}

\begin{defn}[Fórmula]
\ollabel{defn:formulas}
O conjunto~$\Frm[L_0]$ das \emph{!!{formula}s} da lógica proposicional
é definido indutivamente da seguinte maneira:
\begin{enumerate}
\tagitem{prvFalse}{$\lfalse$ é uma !!{formula} atômica.}{}

\tagitem{prvTrue}{$\ltrue$ é uma !!{formula} atômica.}{}

\item Toda !!{propositional variable}~$\Obj p_i$ é uma !!{formula}
  atômica.

\tagitem{prvNot}{Se $!A$ é !!a{formula}, então $\lnot !A$ é
  !!a{formula}.}{}

\tagitem{prvAnd}{Se $!A$ e $!B$ são !!{formula}s, então $(!A \land
  !B)$ é !!a{formula}.}{}

\tagitem{prvOr}{Se $!A$ e $!B$ são !!{formula}s, então $(!A \lor !B)$
  é !!a{formula}.}{}

\tagitem{prvIf}{Se $!A$ e $!B$ são !!{formula}s, então $(!A \lif !B)$
  é !!a{formula}.}{}

\tagitem{prvIff}{Se $!A$ e $!B$ são !!{formula}s, então $(!A \liff !B)$
  é !!a{formula}.}{}

\tagitem{limitClause}{Nada mais é !!a{formula}.}{}
\end{enumerate}
\end{defn}

\begin{explain}
A definição das !!{formula}s é uma \emph{definição indutiva}.
Essencialmente, construímos o conjunto das !!{formula}s em uma
infinidade de estágios. No estágio inicial, declaramos que todas as
fórmulas atômicas são fórmulas; isso corresponde aos primeiros casos
da definição, isto é, aos casos de
\iftag{prvTrue}{$\ltrue$, }{}%
\iftag{prvFalse}{$\lfalse$, }{}%
$\Obj p_i$. Assim, ``!!{formula} atômica'' significa qualquer
!!{formula} dessa forma.

Os outros casos da definição fornecem regras para construir novas
!!{formula}s a partir de !!{formula}s já construídas. No segundo
estágio, podemos usá-las para construir !!{formula}s a partir de
!!{formula}s atômicas. No terceiro estágio, construímos novas fórmulas
a partir das fórmulas atômicas e das obtidas no segundo estágio, e
assim por diante. Uma !!{formula} é qualquer coisa que acaba sendo
construída em algum desses estágios, e nada mais.
\end{explain}

Ao escrever uma fórmula $(!B \ast !C)$ construída a partir de $!B$,
$!C$ mediante um conectivo binário~$\ast$, frequentemente omitiremos
o par exterior de parênteses e escreveremos simplesmente~$!B \ast !C$.

\begin{tagblock}{defNot,defOr,defAnd,defIf,defIff,defTrue,defFalse,defEx,defAll}
\begin{defn}
As fórmulas construídas com os operadores definidos devem ser
entendidas da seguinte maneira:

\begin{tagenumerate}{defTrue,defFalse,defNot,defOr,defAnd,defIf,defIff,defEx,defAll}

\tagitem{defTrue}{$\ltrue$ abrevia
  \iftag{prvFalse}{$\lnot\lfalse$}{$(!A \lor \lnot !A)$ para alguma
    !!{formula} atômica fixa~$!A$}.}{}

\tagitem{defFalse}{$\lfalse$ abrevia
  \iftag{prvTrue}{$\lnot\ltrue$}{$(!A \land \lnot !A)$ para alguma
    !!{formula} atômica fixa~$!A$}.}{}

\tagitem{defNot}{$\lnot !A$ abrevia $!A \lif \lfalse$.}{}

\tagitem{defOr}{$!A \lor !B$ abrevia
  \iftag{prvAnd}{$\lnot(\lnot !A \land \lnot !B)$}{$\lnot !A \lif
    !B$}.}{}

\tagitem{defAnd}{$!A \land !B$ abrevia
  \iftag{prvOr}{$\lnot(\lnot !A \lor \lnot !B)$}{$\lnot (!A \lif
    \lnot !B)$}.}{}

\tagitem{defIf}{$!A \lif !B$ abrevia
  \iftag{prvOr}{$\lnot !A \lor !B)$}{$\lnot (!A \land \lnot !B)$}.}{}

\tagitem{defIff}{$!A \liff !B$ abrevia $(!A \lif !B) \land (!B
  \lif !A)$.}{}
\end{tagenumerate}
\end{defn}
\end{tagblock}

\begin{defn}[Identidade sintática]
O símbolo $\ident$ expressa a identidade sintática entre cadeias de
símbolos, isto é, $!A \ident !B$ sse $!A$ e $!B$ são cadeias de
símbolos de mesmo comprimento que contêm o mesmo símbolo em cada
posição.
\end{defn}

O símbolo $\ident$ pode ser ladeado por cadeias obtidas por
concatenação. Por exemplo, $!A \ident (!B \lor !C)$ significa que a
cadeia de símbolos~$!A$ é a mesma cadeia obtida concatenando, nessa
ordem, um parêntese de abertura, a cadeia $!B$, o símbolo $\lor$, a
cadeia~$!C$ e um parêntese de fechamento. Nesse caso, sabemos que o
primeiro símbolo de $!A$ é um parêntese de abertura, que $!A$ contém
$!B$ como subcadeia (a partir do segundo símbolo), que essa subcadeia
é seguida por $\lor$ etc.

\end{document}
